\chapter*{Introduction G\'en\'erale}
\label{chap:introduction_generale}
\addcontentsline{toc}{chapter}{Introduction G\'en\'erale}

Dans un contexte \'economique marqu\'e par une acc\'el\'eration de la transformation num\'erique, les entreprises de services et d'int\'egration logicielle font face \`a des exigences croissantes en mati\`ere de fiabilit\'e, de maintenabilit\'e et d'exp\'erience utilisateur. Les outils m\'etier internes, longtemps consid\'er\'es comme des projets secondaires, sont devenus des leviers strat\'egiques d'efficacit\'e op\'erationnelle, dont la qualit\'e conditionne directement la productivit\'e des \'equipes qui les utilisent au quotidien.

C'est dans ce contexte que s'inscrit le stage r\'ealis\'e au sein de \textbf{VISEO}, soci\'et\'e fran\c{c}aise de conseil et d'int\'egration num\'erique, sur le projet \textbf{Bifr\"ost}~: une plateforme web de gestion de devis et d'abonnements d\'edi\'ee aux agents du service client. Bifr\"ost repose sur une architecture h\'et\'erog\`ene combinant un backend Python~/ Django REST Framework, une API .NET ASP.NET Core et un frontend Vue.js~/ TypeScript -- une combinaison qui offre flexibilit\'e et robustesse, mais impose une discipline d'\'equipe rigoureuse en mati\`ere de tests, de CI/CD et de qualit\'e de code.

La probl\'ematique centrale de ce projet r\'eside dans la capacit\'e \`a contribuer efficacement \`a un produit logiciel en production~: corriger des anomalies sans introduire de r\'egressions, implanter de nouvelles fonctionnalit\'es dans le respect de l'architecture existante, et am\'eliorer progressivement la qualit\'e du code sans d\'estabiliser les flux en cours d'utilisation. Cette contrainte, propre aux projets en TMA (Tierce Maintenance Applicative), a guid\'e l'ensemble des d\'ecisions techniques tout au long du stage.

Le pr\'esent rapport s'articule autour de quatre chapitres principaux. Le \textbf{premier chapitre} pr\'esente l'entreprise d'accueil VISEO, son positionnement et le contexte dans lequel s'inscrit le projet Bifr\"ost. Le \textbf{deuxi\`eme chapitre} analyse l'existant technique~: architecture, stack, organisation des \'equipes et processus de d\'eveloppement. Le \textbf{troisi\`eme chapitre} d\'etaille les r\'ealisations techniques effectu\'ees sur les trois couches de l'application (frontend, API .NET, backend Django), en explicitant les choix d'impl\'ementation et les contraintes rencontr\'ees. Enfin, le \textbf{quatri\`eme chapitre} aborde les perspectives d'\'evolution, notamment l'initiative d'observabilit\'e et la strat\'egie de mise \`a jour des d\'ependances Python.